%\title{University of Bristol Thesis Template}
\RequirePackage[l2tabu]{nag}		% Warns for incorrect (obsolete) LaTeX usage
%
%
% File: memoirthesis.tex
% Author: Victor Brena
% Description: Contains the thesis template using memoir class,
% which is mainly based on book class but permits better control of 
% chapter styles for example. This template is an adaptation and 
% modification of Oscar's.
% 
% Memoir is a flexible class for typesetting poetry, fiction, 
% non-fiction and mathematical works as books, reports, articles or
% manuscripts. CTAN repository is found at:
% http://www.ctan.org/tex-archive/macros/latex/contrib/memoir/
%
%
% UoB guidelines for thesis presentation were found at:
% http://www.bris.ac.uk/esu/pg/pgrcop11-12topic.pdf#page=49
%
% UoB guidelines:
%
% The dissertation must be printed on A4 white paper. Paper up to A3 may be used
% for maps, plans, diagrams and illustrative material. Pages (apart from the
% preliminary pages) should normally be double-sided.
%
% Memoir class loads useful packages by default (see manual).
\documentclass[a4paper,11pt,leqno,openbib,oldfontcommands]{memoir} %add 'draft' to turn draft option on (see below)


%
% Adding metadata:
\usepackage{datetime}
\usepackage{ifpdf}
\ifpdf
\pdfinfo{
   /Author (Fergus J. Moore)
   /Title (PhD Thesis)
   /Keywords (One; Two;Three)
   /CreationDate (D:\pdfdate)
}
\fi
% When draft option is on. 
\ifdraftdoc 
	\usepackage{draftwatermark}				%Sets watermarks up.
	\SetWatermarkScale{0.3}
	\SetWatermarkText{\bf Draft: \today}
\fi
%
% Declare figure/table as a subfloat.
\newsubfloat{figure}
\newsubfloat{table}
% Better page layout for A4 paper, see memoir manual.
\settrimmedsize{297mm}{210mm}{*}
\setlength{\trimtop}{0pt} 
\setlength{\trimedge}{\stockwidth} 
\addtolength{\trimedge}{-\paperwidth} 
\settypeblocksize{634pt}{448.13pt}{*} 
\setulmargins{4cm}{*}{*} 
\setlrmargins{*}{*}{1.5} 
\setmarginnotes{17pt}{51pt}{\onelineskip} 
\setheadfoot{\onelineskip}{2\onelineskip} 
\setheaderspaces{*}{2\onelineskip}{*} 
\checkandfixthelayout
%
\frenchspacing
% Font with math support: New Century Schoolbook
\usepackage{fouriernc}
\usepackage[T1]{fontenc}
%
% UoB guidelines:
%
% Text should be in double or 1.5 line spacing, and font size should be
% chosen to ensure clarity and legibility for the main text and for any
% quotations and footnotes. Margins should allow for eventual hard binding.
%
% Note: This is automatically set by memoir class. Nevertheless \OnehalfSpacing 
% enables double spacing but leaves single spaced for captions for instance. 
%\DoubleSpacing
%\OnehalfSpacing 
%
% Sets numbering division level
\setsecnumdepth{subsection} 
\maxsecnumdepth{subsubsection}
%
% Chapter style (taken and slightly modified from Lars Madsen Memoir Chapter 
% Styles document
\usepackage[dvipsnames]{xcolor}
\usepackage{calc,soul,fourier}
\makeatletter 
\newlength\dlf@normtxtw 
\setlength\dlf@normtxtw{\textwidth} 
\newsavebox{\feline@chapter} 
\newcommand\feline@chapter@marker[1][4cm]{%
	\sbox\feline@chapter{% 
		\resizebox{!}{#1}{\fboxsep=1pt%
			\colorbox{Tan}{\color{white}\thechapter}% 
		}}%
		\rotatebox{90}{% 
			\resizebox{%
				\heightof{\usebox{\feline@chapter}}+\depthof{\usebox{\feline@chapter}}}% 
			{!}{\scshape\so\@chapapp}}\quad%
		\raisebox{\depthof{\usebox{\feline@chapter}}}{\usebox{\feline@chapter}}%
} 
\newcommand\feline@chm[1][4cm]{%
	\sbox\feline@chapter{\feline@chapter@marker[#1]}% 
	\makebox[0pt][c]{% aka \rlap
		\makebox[1cm][r]{\usebox\feline@chapter}%
	}}
\makechapterstyle{daleifmodif}{
	\renewcommand\chapnamefont{\normalfont\Large\scshape\raggedleft\so} 
	\renewcommand\chaptitlefont{\normalfont\Large\bfseries\scshape} 
	\renewcommand\chapternamenum{} \renewcommand\printchaptername{} 
	\renewcommand\printchapternum{\null\hfill\feline@chm[2.5cm]\par} 
	\renewcommand\afterchapternum{\par\vskip\midchapskip} 
	\renewcommand\printchaptertitle[1]{\color{gray}\chaptitlefont\raggedleft ##1\par}
} 
\makeatother 
\chapterstyle{daleifmodif}
%
% UoB guidelines:
%
% The pages should be numbered consecutively at the bottom centre of the
% page.
\makepagestyle{myvf} 
\makeoddfoot{myvf}{}{\thepage}{} 
\makeevenfoot{myvf}{}{\thepage}{} 
\makeheadrule{myvf}{\textwidth}{\normalrulethickness} 
\makeevenhead{myvf}{\small\textsc{\leftmark}}{}{} 
\makeoddhead{myvf}{}{}{\small\textsc{\rightmark}}
\pagestyle{myvf}
%
% Oscar's command (it works):
% Fills blank pages until next odd-numbered page. Used to emulate single-sided
% frontmatter. This will work for title, abstract and declaration. Though the
% contents sections will each start on an odd-numbered page they will
% spill over onto the even-numbered pages if extending beyond one page
% (hopefully, this is ok).
\newcommand{\clearemptydoublepage}{\newpage{\thispagestyle{empty}\cleardoublepage}}
%
%
% Creates indexes for Table of Contents, List of Figures, List of Tables and Index
\makeindex
% \printglossaries below creates a list of abbreviations. \gls and related
% commands are then used throughout the text, so that latex can automatically
% keep track of which abbreviations have already been defined in the text.
%
% The import command enables each chapter tex file to use relative paths when
% accessing supplementary files. For example, to include
% chapters/brewing/images/figure1.png from chapters/brewing/brewing.tex we can
% use
% \includegraphics{images/figure1}
% instead of
% \includegraphics{chapters/brewing/images/figure1}
\usepackage{import}

% Add other packages needed for chapters here. For example:
\usepackage{lipsum}					%Needed to create dummy text
\usepackage{amsfonts} 					%Calls Amer. Math. Soc. (AMS) fonts
\usepackage[centertags]{amsmath}			%Writes maths centred down
\usepackage{stmaryrd}					%New AMS symbols
\usepackage{amssymb}					%Calls AMS symbols
\usepackage{amsthm}					%Calls AMS theorem environment
\usepackage{newlfont}					%Helpful package for fonts and symbols
\usepackage{layouts}					%Layout diagrams
\usepackage{graphicx}					%Calls figure environment
\usepackage{longtable,rotating}			%Long tab environments including rotation. 
\usepackage[utf8]{inputenc}			%Needed to encode non-english characters 
									%directly for mac
\usepackage{colortbl}					%Makes coloured tables
\usepackage{wasysym}					%More math symbols
\usepackage{mathrsfs}					%Even more math symbols
\usepackage{float}						%Helps to place figures, tables, etc. 
\usepackage{verbatim}					%Permits pre-formated text insertion
\usepackage{upgreek }					%Calls other kind of greek alphabet
\usepackage{latexsym}					%Extra symbols
\usepackage[square,numbers,
		     sort&compress]{natbib}		%Calls bibliography commands 
\usepackage{url}						%Supports url commands
% \usepackage{etex}						%eTeXÕs extended support for counters
% \usepackage{fixltx2e}					%Eliminates some in felicities of the 
									%original LaTeX kernel
\usepackage[spanish,english]{babel}		%For languages characters and hyphenation
\usepackage{color}                    				%Creates coloured text and background
\usepackage[colorlinks=true,
		     allcolors=black]{hyperref}              %Creates hyperlinks in cross references
\usepackage{memhfixc}					%Must be used on memoir document 
									%class after hyperref
\usepackage{enumerate}					%For enumeration counter
\usepackage{footnote}					%For footnotes
\usepackage{microtype}					%Makes pdf look better.
\usepackage{rotfloat}					%For rotating and float environments as tables, 
									%figures, etc. 
\usepackage{alltt}						%LaTeX commands are not disabled in 
									%verbatim-like environment
\usepackage[version=0.96]{pgf}			%PGF/TikZ is a tandem of languages for producing vector graphics from a 
\usepackage{tikz}						%geometric/algebraic description.
\usetikzlibrary{arrows,shapes,snakes,
		       automata,backgrounds,
		       petri,topaths}				%To use diverse features from tikz		
%							
%Reduce widows  (the last line of a paragraph at the start of a page) and orphans 
% (the first line of paragraph at the end of a page)
\widowpenalty=1000
\clubpenalty=1000
%
%
% My caption style
\newcommand{\mycaption}[2][\@empty]{
	\captionnamefont{\scshape} 
	\changecaptionwidth
	\captionwidth{0.9\linewidth}
	\captiondelim{.\:} 
	\indentcaption{0.75cm}
	\captionstyle[\centering]{}
	\setlength{\belowcaptionskip}{10pt}
	\ifx \@empty#1 \caption{#2}\else \caption[#1]{#2}
}
%
% My subcaption style
\newcommand{\mysubcaption}[2][\@empty]{
	\subcaptionsize{\small}
	\hangsubcaption
	\subcaptionlabelfont{\rmfamily}
	\sidecapstyle{\raggedright}
	\setlength{\belowcaptionskip}{10pt}
	\ifx \@empty#1 \subcaption{#2}\else \subcaption[#1]{#2}
}
%
%An initial of the very first character of the content
\usepackage{lettrine}
\newcommand{\initial}[1]{%
	\lettrine[lines=3,lhang=0.33,nindent=0em]{
		\color{Tan}
     		{\textsc{#1}}}{}}
%
%
\hyphenation{res-pec-tively}
\hyphenation{mono-ti-ca-lly}
\hyphenation{hypo-the-sis}
\hyphenation{para-me-ters}
\hyphenation{sol-va-bi-li-ty}
%
%

\newcommand{\pecl}{\operatorname{\mathit{P\kern-.08em e}}}
\newcommand{\fergus}[1]{{\color{orange} #1}}
\newcommand{\latin}[1]{{\itshape #1}}
\newcommand{\etal}{\latin{et al.}}

\usepackage{todonotes}
\usepackage{enumitem}  

\hypersetup{colorlinks=true, linkcolor=blue, urlcolor=blue} 


\begin{document}

% UoB guidlines:
%
% Preliminary pages
% 
% The five preliminary pages must be the Title Page, Abstract, Dedication
% and Acknowledgements, Author's Declaration and Table of Contents.
% These should be single-sided.
% 
% Table of contents, list of tables and illustrative material
% 
% The table of contents must list, with page numbers, all chapters,
 % sections and subsections, the list of references, bibliography, list of
% abbreviations and appendices. The list of tables and illustrations
% should follow the table of contents, listing with page numbers the
% tables, photographs, diagrams, etc., in the order in which they appear
% in the text.
% 
\frontmatter
\pagenumbering{roman}
%
%
% File: Title.tex
% Author: V?ctor Bre?a-Medina
% Description: Contains the title page
%
% UoB guidelines:
% 
% At the top of the title page, within the margins, the dissertation should give the title and, if 
% necessary, sub-title and volume number. If the dissertation is in a language other than English, the 
% title must be given in that language and in English. The full name of the author should be in the 
% centre of the page. At the bottom centre should be the words ?A dissertation submitted to the 
% University of Bristol in accordance with the requirements for award of the degree of ? in the 
% Faculty of ...?, with the name of the school and month and year of submission. The word count of 
% the dissertation (text only) should be entered at the bottom right-hand side of the page.
%
%
\begin{titlingpage}
\begin{SingleSpace}
\calccentering{\unitlength} 
\begin{adjustwidth*}{\unitlength}{-\unitlength}
\vspace*{13mm}
\begin{center}
\rule[0.5ex]{\linewidth}{2pt}\vspace*{-\baselineskip}\vspace*{3.2pt}
\rule[0.5ex]{\linewidth}{1pt}\\[\baselineskip]
{\HUGE Crystallisation and Confinement of Mesoscale Active Matter }\\[4mm]
%{\Large \textit{Subtitle}}\\
\rule[0.5ex]{\linewidth}{1pt}\vspace*{-\baselineskip}\vspace{3.2pt}
\rule[0.5ex]{\linewidth}{2pt}\\
\vspace{6.5mm}
{\large By}\\
\vspace{6.5mm}
{\large\textsc{Fergus J. Moore}}\\
\vspace{11mm}
\includegraphics[scale=0.6]{logos/bristolcrest_colour}\\
\vspace{6mm}
{\large Department of Physics\\
\textsc{University of Bristol}}\\
\vspace{11mm}
\begin{minipage}{10cm}
A dissertation submitted to the University of Bristol in accordance with the requirements of the degree of \textsc{Doctor of Philosophy} in the Faculty of Physics.
\end{minipage}\\
\vspace{9mm}
{\large\textsc{\today}}
\vspace{12mm}
\end{center}
\begin{flushright}
{\small Word count: 32152 }
\end{flushright}
\end{adjustwidth*}
\end{SingleSpace}
\end{titlingpage}
\clearemptydoublepage
%
%
% File: abstract.tex
% Author: V?ctor Bre?a-Medina
% Description: Contains the text for thesis abstract
%
% UoB guidelines:
%
% Each copy must include an abstract or summary of the dissertation in not
% more than 300 words, on one side of A4, which should be single-spaced in a
% font size in the range 10 to 12. If the dissertation is in a language other
% than English, an abstract in that language and an abstract in English must
% be included.

\chapter*{Abstract}
\begin{SingleSpace}

\initial{A}ctive matter systems are comprised of bodies that consume energy from the environment for self-propulsion and as a result are driven out-of-equilibrium. These systems encompass a range of lengthscales and describe both living and artificial systems. Moreover, the motion and collective interactions of these active bodies lead to an ensemble of rich and complex behaviours not seen in equilibrium systems. On mesoscopic lengthscales (nm to $\mu$m), the prominent active systems are bacterial swimmers and active colloidal particles.
 However, as a result of the experimental challenge involved, there have been relatively few studies of mesoscale active matter systems in 3D. In this thesis, we aim to address two poorly understood areas of active systems in 3D: the crystallisation of active Brownian spheres, and the behaviour of active matter in 3D non-trivial confinement. 

Regarding crystallisation, we perform molecular dynamics simulations to explore nucleation and polymorph selection in crystals of mono-disperse active spheres. We find two activity dependant crystallisation regimes: spinodal, and nucleation and growth. In the former we observe an increase in the crystal growth rate due to the accelerated dynamics in the melt, and in the latter we find that activity inhibits nucleation. Remarkably, activity also has a strong effect on polymorph selection. While the passive system crystallises in an equimolar mixture of FCC and HCP, active particles progressively favour the FCC phase at higher activity.

The transport of mesoscale active matter occurs naturally in complex environments that are often crowded, random and irregular. We explore structural and dynamical properties of 3D active particles in random environments with a multidisciplinary approach. First we conduct molecular dynamics simulations of active Brownian spheres in two random environments. We find that the confinement increases the heterogeneity of the dynamics, with new populations of absorbed and localised particles appearing close to the obstacles. This heterogeneity has a profound impact on the motility induced phase separation  exhibited by the particles at high activity, ranging from growth in random disorder to a bicontinuous phase separation in porous environments.

Finally we approach this same problem through experiment and present the first 3D active colloidal system in non-trivial confinement. The system is comprised of active Janus particles driven via an external field, propelling through the porous network surrounding a colloidal gel. 
 This system is imaged in real-time with confocal microscopy and the particle dynamics are extracted with particle tracking methods. 



\end{SingleSpace}
\clearpage
\clearemptydoublepage
%

\chapter*{Acknowledgements}
\begin{SingleSpace}


\initial{T}his work would not exist without the support of a great many people. I would first like to thank John Russo and Paddy Royall, without whom this project would not exist. Their advice and support has been foundational to my journey from student to scientist, and has remained a constant despite their departures from Bristol. I would also like to thank Tanniemola Liverpool and Henkjan Gersen, who adopted me late into my project and gave the advice and encouragement I needed to get this work over the line.

I consider myself incredibly fortunate to have undertaken my PhD through the BCFN. No one could wish for a warmer welcome than I received from the community cultivated by Annela Seddon and Ian Lindsey. The word \textit{family} has often been used internally to describe the BCFN and it is a sentiment I share enthusiastically. I began my studies through the BCFN in 2017 with a cohort of wonderful people, we have all undergone this journey together and I'm immensely proud of all of you. 

This work was predominantly conducted in the company of the students and postdocs who make up the Royall group. In many ways this work stands upon the small day to day interactions taking place in the lab, the office, or the pub. I have become very fond of you all, and I would especially like to thank Abraham, Yushi, Katherine, Nariaki, Josh, Francesco and Ioatzin, for many helpful discussions that have been essential to my getting here. I would also like to thank Claire Noble, as the other half of the "Biochemistry thesis club" it has been my pleasure to write in your company.

For me, as is the case with many people, I have found that there is an intangible part to a PhD that leaves the university and settles into the corner of whichever room one finds themselves in.
In regard to this aspect I would like to give thanks to a number of people outside of academia who have made this work possible: the various inhabitants and associates of 15d Beaufort road, especially Sam Hueget, Sam Smithers, Will Dixon and Peter Rosso, who's real and imaginary parts have performed great feats of heroism over these past years. Teo Garcia Millan, and many others who have helped me to climb to new heights. Flynn, who is cute and fluffy. My parents, who are my greatest champions. Their love and support have made innumerable contributions, \newline I love you both. Finally, Rachel, you have been here through all of the highs and lows, for the nights when reaching this moment seemed impossible, and popping prosecco at calls for eureka. I am forever grateful for your support. I love you.     



\end{SingleSpace}
\clearpage
\clearemptydoublepage
%
%
% File: declaration.tex
% Author: V?ctor Bre?a-Medina
% Description: Contains the declaration page
%
% UoB guidelines:
%
% Author's declaration
%
% I declare that the work in this dissertation was carried out in accordance
% with the requirements of the University's Regulations and Code of Practice
% for Research Degree Programmes and that it has not been submitted for any
% other academic award. Except where indicated by specific reference in the
% text, the work is the candidate's own work. Work done in collaboration with,
% or with the assistance of, others, is indicated as such. Any views expressed
% in the dissertation are those of the author.
%
% SIGNED: .............................................................
% DATE:..........................
%
\chapter*{Author's declaration}
\begin{SingleSpace}
\begin{quote}
\initial{I} declare that the work in this dissertation was carried out in accordance with the requirements of  the University's Regulations and Code of Practice for Research Degree Programmes and that it  has not been submitted for any other academic award. Except where indicated by specific  reference in the text, the work is the candidate's own work. Work done in collaboration with, or with the assistance of, others, is indicated as such. Any views expressed in the dissertation are those of the author.

\vspace{1.5cm}
\noindent
\hspace{-0.75cm}\textsc{SIGNED: .................................................... DATE: ..........................................}
\end{quote}
\end{SingleSpace}
\clearpage



\section*{Previous Publications}


A section of this thesis has appeared in the following paper:


\begin{enumerate}
	\item \textbf{Chapter \ref{chap:crystal}}:   F. J. Moore, T. Liverpool, C.P. Royall and J. Russo, "Crystallisation and Polymorph Selection in Active Brownian Particles", \textit{European Physical Journal - E}, 44 (9), 1-8 (2021).
\end{enumerate}

\noindent In addition, the author contributed to a number of other publications while working towards this qualification:

\begin{enumerate}
	\item J. F. Robinson, I. Rios de Anda, F. J. Moore, J. P. Reid, R. P. Sear, and C. P. Royall, "Efficacy of face coverings in reducing transmission of COVID-19: Calculations based on models of droplet capture." \textit{Physics of fluids},  33 (4), 043112 (2021).
	\item J. F. Robinson, I. Rios de Anda, F. J. Moore, F. K. Gregson, J. P. Reid, L. Husain, R. P. Sear, and C. P. Royall, "How effective are face coverings in reducing transmission of COVID-19?" \textit{Aerosol Science and Technology}, 1-29  (2022).
\end{enumerate}
\clearemptydoublepage
%
\renewcommand{\contentsname}{Table of Contents}
\maxtocdepth{subsection}
\tableofcontents*
\addtocontents{toc}{\par\nobreak \mbox{}\hfill{\bf Page}\par\nobreak}
\clearemptydoublepage
%




%\listoftables
%\addtocontents{lot}{\par\nobreak\textbf{{\scshape Table} \hfill Page}\par\nobreak}
%\clearemptydoublepage


%
\listoffigures
\addtocontents{lof}{\par\nobreak\textbf{{\scshape Figure} \hfill Page}\par\nobreak}
\clearemptydoublepage
%
%
% The bulk of the document is delegated to these chapter files in
% subdirectories.

\DoubleSpacing

\mainmatter



%
%\import{chapters/introduction/}{introduction.tex}
\import{chapters/activeMatter/}{activeMatter.tex}
\import{chapters/colloids/}{colloids.tex}
\import{chapters/expMethods/}{expMethods.tex}
\import{chapters/aprilFools/}{aprilFools.tex}
\import{chapters/activeConfinement/}{activeConfinement.tex}
\import{chapters/expSystem/}{expSystem.tex}
\import{chapters/conclusion/}{conclusion.tex}

\clearemptydoublepage
%
%
% And the appendix goes here
\appendix
\import{chapters/appendices/}{app0A.tex}
\clearemptydoublepage
%
% Apparently the guidelines don't say anything about citations or
% bibliography styles so I guess we can use anything.
\backmatter
\bibliographystyle{siam}
\refstepcounter{chapter}
\bibliography{fergusThesis.bib}
\clearemptydoublepage
%
% Add index
%\printindex
%   
\end{document}